\documentclass[conference,compsoc]{IEEEtran}

\usepackage{amsmath,amsfonts,amssymb}
\usepackage{algorithmic}
\usepackage{algorithm}
\usepackage{array}
\usepackage[caption=false,font=footnotesize,labelfont=sf,textfont=sf]{subfig}
\usepackage{textcomp}
\usepackage{stfloats}
\usepackage{url}
\usepackage{verbatim}
\usepackage{graphicx}
\usepackage[nocompress]{cite}
\usepackage{listings}
\usepackage{xcolor}
\usepackage[catalan]{babel}
\usepackage[T1]{fontenc}
\usepackage[utf8]{inputenc}
\usepackage{booktabs}
\usepackage{float}
\usepackage{cuted}
\usepackage{pgfplots}
\pgfplotsset{compat=1.18}
\graphicspath{{./}}

\renewcommand{\abstractname}{Resum}
\renewcommand{\IEEEkeywordsname}{Termes d'índex}

\makeatletter
\def\abstract{\normalfont\@IEEEtweakunitybaselinestretch{1.15}%
    \if@twocolumn
      \@IEEEabskeysecsize\noindent\textbf{\textit{\abstractname}}---\normalfont\@IEEEabskeysecsize\relax
    \else
      \bgroup\par\addvspace{0.5\baselineskip}\centering\vspace{-1.78ex}\@IEEEabskeysecsize\textbf{\abstractname}\par\addvspace{0.5\baselineskip}\egroup\quotation\normalfont\@IEEEabskeysecsize%
    \fi\@IEEEgobbleleadPARNLSP}
\def\IEEEkeywords{\normalfont\@IEEEtweakunitybaselinestretch{1.15}%
    \if@twocolumn
      \@IEEEabskeysecsize\vskip 0.5\baselineskip plus 0.25\baselineskip minus 0.25\baselineskip\noindent
      \textbf{\textit{\IEEEkeywordsname}}---\normalfont\@IEEEabskeysecsize\relax
    \else
      \bgroup\par\addvspace{0.5\baselineskip}\centering\@IEEEabskeysecsize\textbf{\IEEEkeywordsname}\par\addvspace{0.5\baselineskip}\egroup\quotation\normalfont\@IEEEabskeysecsize%
    \fi\@IEEEgobbleleadPARNLSP}
\makeatother

\lstset{
    language=Java,
    basicstyle=\ttfamily\footnotesize,
    keywordstyle=\color{blue}\bfseries,
    commentstyle=\color{gray}\itshape,
    stringstyle=\color{red!70!black},
    numbers=left,
    numberstyle=\tiny\color{gray},
    numbersep=5pt,
    frame=single,
    breaklines=true,
    showstringspaces=false,
    tabsize=4,
    captionpos=b
}

\hyphenation{pro-gra-ma-cio di-na-mi-ca me-mo-it-za-cio cal-cu-la-do-ra con-tro-la-dor}

\begin{document}

\title{Joc de la Calculadora amb Programació Dinàmica\\Pràctica 5}

\author{\IEEEauthorblockN{Serafí Nebot Ginard}
\IEEEauthorblockA{Universitat de les Illes Balears (UIB)\\
Grau en Enginyeria Informàtica\\
Algorismes Avançats, Curs 2025--2026}}

\maketitle

\begin{strip}
\begin{abstract}
Aquesta memòria descriu una aplicació Java que resol el joc de la calculadora
mitjançant programació dinàmica. El joc parteix d'un total acumulat, una tecla
jugada anteriorment i un jugador a qui toca moure; a cada torn només es poden
prémer tecles de la mateixa fila o columna que la tecla anterior, i perd el
jugador que arriba o supera el límit. S'han implementat dues variants de la
mateixa recurrència: una versió \textit{top-down} amb memoització i una versió
\textit{bottom-up} que precomputa tota la taula d'estats. L'aplicació separa
model, vista i controlador, ofereix una interfície Swing amb recàlcul automàtic,
teclats estàndard o aleatoris i una reproducció pas a pas de la partida òptima.
També s'ha incorporat una eina de benchmark que mesura temps de CPU, temps real,
estats calculats i variació aproximada de memòria. Els resultats confirmen que
la versió \textit{bottom-up} té un cost més predictible i proporcional a tota la
taula, mentre que la versió \textit{top-down} visita menys estats quan la partida
inicial només arriba a una part de l'espai possible.
\end{abstract}

\begin{IEEEkeywords}
programació dinàmica, jocs combinatoris, memoització, taula \textit{bottom-up},
Model-Vista-Controlador, Java Swing, benchmark
\end{IEEEkeywords}
\end{strip}

\IEEEpeerreviewmaketitle

%% ================================================================
\section{Introducció}

La pràctica tracta un joc determinista sobre una calculadora. Hi ha un valor
acumulat, un límit de pèrdua, un teclat rectangular i un conjunt de jugadors.
En cada torn, el jugador prem una tecla per sumar el seu valor al total. Si el
nou total arriba o supera el límit, aquest jugador perd immediatament. La
restricció que fa interessant el problema és que, després de la primera jugada,
només es poden emprar tecles situades a la mateixa fila o columna que la tecla
jugada anteriorment.

L'objectiu principal és determinar quin jugador perd si tots juguen de manera
òptima per evitar ser el perdedor. En partides de dos jugadors això també
identifica el guanyador; en partides amb més jugadors només es reporta el
perdedor, perquè hi pot haver més d'un supervivent i no hi ha un únic guanyador
natural. La noció d'òptim emprada és local a cada jugador: en el seu torn, cada
jugador tria un moviment que evita la seva pròpia derrota si existeix. No es
modela cooperació entre jugadors ni preferències sobre quin altre jugador hauria
de perdre.

La solució implementada no explora l'arbre de joc de manera ingènua. Com que el
total acumulat sempre augmenta, el graf d'estats és acíclic i es pot resoldre
amb programació dinàmica \cite{ref_gfg_dp}. La pràctica compara dues formes de
la mateixa idea: una a demanda amb memoització i una iterativa que omple tota la
taula. A més, s'ha construït una aplicació d'escriptori completa amb patró
Model-Vista-Controlador i una eina separada per obtenir dades experimentals.

%% ================================================================
\section{Descripció del Problema}

El teclat té amplada $W$ i alçada $H$, amb dimensions entre 2 i 5. El nombre de
tecles és $K = W \cdot H$ i cada tecla conté exactament un valor de l'interval
$1,\ldots,K$. La distribució estàndard imita una calculadora: per exemple, un
teclat $3\times3$ mostra $7,8,9$ a la fila superior, $4,5,6$ al centre i
$1,2,3$ a la part inferior.

El valor especial $0$ indica que encara no s'ha premut cap tecla. En aquest cas,
totes les tecles són moviments vàlids. Si la darrera tecla jugada és $l>0$, el
conjunt de moviments vàlids $M(l)$ conté les tecles de la mateixa fila i de la
mateixa columna que $l$, excloent $l$ mateix. Per exemple, en el teclat
$3\times3$ estàndard, si la darrera tecla és $6$, els moviments vàlids són
$3,4,5,9$.

La configuració d'una partida queda definida per:

\begin{itemize}
    \item el teclat concret;
    \item el total inicial $t_0$;
    \item el límit de pèrdua $L$;
    \item la darrera tecla jugada $l_0$, o $0$ si no n'hi ha;
    \item el nombre de jugadors $P$;
    \item el jugador que ha de moure.
\end{itemize}

La implementació permet teclats estàndard i teclats aleatoris. En mode aleatori,
la disposició continua essent una permutació vàlida de $1,\ldots,K$; només canvia
la posició de cada valor. Això obliga la generació aleatòria a preservar dues
invariants: no hi pot haver valors repetits i no pot faltar cap tecla de
l'interval. A la interfície, el botó \textit{Nou aleatori} canvia la llavor usada
per generar la permutació, mentre que el benchmark empra teclats estàndard per
fer comparables les mesures.

%% ================================================================
\section{Fonaments Teòrics}

\subsection{Estat i Recurrència}

L'estat resolt pels algorismes és:

\begin{equation}
S = (t, l, j)
\end{equation}

on $t$ és el total actual, $l$ és la darrera tecla jugada i $j$ és el jugador a
qui toca moure. La funció dinàmica $F(t,l,j)$ retorna l'índex del jugador que
perdrà des d'aquest estat si tots juguen òptimament.

Per a cada moviment vàlid $m \in M(l)$ es defineix el successor segur:

\begin{equation}
\operatorname{succ}(t,l,j,m) = (t+m, m, (j+1) \bmod P)
\end{equation}

sempre que $t+m<L$. Si $t+m \geq L$, el moviment és terminal i fa perdre
directament el jugador $j$, de manera que no és un moviment útil per evitar la
pròpia derrota.
Per abreujar la notació, escrivim $j'=(j+1)\bmod P$.

El conjunt de moviments que salven el jugador actual és:

\begin{multline}
A(t,l,j)=\{m\in M(l)\mid t+m<L\\
{}\land F(t+m,m,j')\neq j\}
\end{multline}

La recurrència implementada és:

\begin{equation}
F(t,l,j)=
\begin{cases}
j,               & \text{si } A(t,l,j) = \emptyset\\
F(t+m^*,m^*,j'), & \text{si } A(t,l,j) \neq \emptyset\\
\end{cases}
\label{eq:recurrence}
\end{equation}

on $m^*$ és el primer moviment de $A(t,l,j)$ en ordre ascendent. Aquest desempat
no canvia si el jugador pot evitar perdre, però fa que la reproducció de la
partida sigui determinista i que els dos algorismes mostrin exactament la mateixa
línia òptima.

\subsection{Per què Programació Dinàmica}

Una exploració directa podria recalcular moltes vegades el mateix estat, perquè
seqüències diferents de tecles poden conduir al mateix triple $(t,l,j)$. La
programació dinàmica evita aquesta repetició guardant el valor de cada estat una
sola vegada.

El detall clau és que no hi ha cicles: qualsevol moviment vàlid incrementa el
total en un valor positiu. Per tant, tots els successors segurs d'un estat tenen
un total més gran. Això permet resoldre el problema tant de forma recursiva amb
memoització (\textit{top-down}) com de forma iterativa baixant de $L-1$ fins a $0$
(\textit{bottom-up}).

En la variant \textit{top-down}, el programa parteix només de l'estat inicial i
aplica la recurrència de manera recursiva. Quan apareix un estat nou, es calcula
el seu perdedor i es guarda en una memòria; si el mateix estat torna a aparèixer
per un altre camí de joc, es reutilitza el valor ja calculat. Per tant, aquesta
variant calcula només la part de l'espai d'estats que és necessària per respondre
la consulta inicial.

En la variant \textit{bottom-up}, en canvi, no es comença per l'estat inicial,
sinó que s'omple tota la taula d'estats possibles. Es processen primer els totals
més propers al límit, perquè en aquests casos els successors ja són terminals o
molt simples. Després es baixa progressivament fins a $0$. Quan es calcula un
estat $(t,l,j)$, tots els successors amb total $t+m$ ja estan resolts, i per això
el valor es pot obtenir consultant directament la taula.

\subsection{Correcció dels Algorismes}

La correcció es basa en inducció sobre la distància fins al límit. Els estats amb
totals propers a $L$ són els més simples: si tots els moviments vàlids arriben o
superen el límit, qualsevol jugada fa perdre el jugador actual i, per tant,
$F(t,l,j)=j$. Aquest és el cas base implícit de la recurrència.

Per al pas inductiu, es considera un estat $(t,l,j)$ amb successors segurs. Com
que tots els moviments sumen un valor positiu, qualsevol successor té total
$t+m>t$ i està més a prop del límit. Per hipòtesi inductiva, el valor
$F(t+m,m,j')$ d'aquests successors ja descriu correctament qui perd des d'allà.
Si existeix qualque moviment segur amb $F(t+m,m,j')\neq j$, el jugador actual pot
evitar ser el perdedor triant-lo. Si no existeix cap moviment d'aquest tipus,
tots els camins possibles acaben fent perdre el jugador actual, i per això la
resposta correcta és $j$.

Les dues variants implementen aquesta mateixa recurrència. La versió
\textit{top-down} la avalua sota demanda i memoïtza els resultats; la versió
\textit{bottom-up} calcula els mateixos valors en un ordre que garanteix que cada
successor segur ja està resolt. Per això, si parteixen de la mateixa configuració,
han de retornar el mateix perdedor i la mateixa primera jugada segons el criteri
de desempat.

\subsection{Complexitat}

El nombre màxim d'estats no terminals és:

\begin{equation}
S_{max}=L\cdot(K+1)\cdot P
\label{eq:states}
\end{equation}

El factor $K+1$ apareix perquè $l$ pot ser qualsevol tecla o el valor especial
$0$. Per estat, el nombre de moviments vàlids és com a màxim:

\begin{equation}
B=\max(K, W+H-2)
\end{equation}

El cas $K$ correspon al primer moviment, quan $l=0$; la resta d'estats tenen com
a màxim $W+H-2$ moviments perquè només es pot triar fila o columna.

La versió \textit{top-down} només expandeix els estats que la recurrència arriba
a necessitar des de la configuració inicial. Si $R$ és aquest nombre d'estats
visitats, el cost és:

\begin{equation}
T_{TD}=O(R\cdot B), \qquad M_{TD}=O(R)
\end{equation}

Cal distingir $R$ de $S_{max}$, que és la mida de l'espai complet d'estats
que es podria construir, mentre que $R$ és només el nombre d'estats que la
recursió arriba a visitar i guardar per a una configuració inicial concreta. Per
tant, sempre es compleix $R\leq S_{max}$. En el pitjor cas, si l'exploració des
de l'estat inicial acaba necessitant pràcticament tota la taula, $R$ pot arribar
a ser igual o molt proper a $S_{max}$; però en molts casos és menor, i aquest és
el principal avantatge pràctic de la variant \textit{top-down}.

La versió \textit{bottom-up} omple tota la taula independentment de si tots els
estats seran consultats des de la partida inicial:

\begin{equation}
T_{BU}=O(S_{max}\cdot B), \qquad
M_{BU}=O(S_{max})
\end{equation}

Encara que la implementació \textit{bottom-up} guarda dues taules tridimensionals,
una per al jugador perdedor i una altra per al moviment triat, cadascuna té mida
proporcional a $S_{max}$. Per tant, la memòria total és $O(2\cdot S_{max})$, que
asimptòticament continua sent $O(S_{max})$.

%% ================================================================
\section{Disseny i Implementació}

\subsection{Arquitectura MVC}

L'aplicació segueix el patró Model-Vista-Controlador \cite{ref_mvc_wikipedia}. El punt
d'entrada \texttt{Main} crea la finestra Swing i després crea el controlador.
La finestra no coneix cap classe concreta de controlador; només emet esdeveniments a
través de la interfície \texttt{GameViewListener}. Això manté la vista separada
de les regles del joc.

La decisió de fer aquesta separació no és només organitzativa. El projecte té
dues parts molt diferents: d'una banda hi ha la solució del problema amb algorismes
de programació dinàmica; de l'altra, una interfície interactiva que ha de respondre
a cada canvi de l'usuari. Si la vista contingués lògica de joc, seria difícil validar
els casos límit sense obrir finestres. Si el model conegués Swing, els algorismes
no es podrien reutilitzar en el benchmark. Per això la dependència sempre va en
la mateixa direcció: la vista exposa una abstracció, el controlador la consumeix
i el model no coneix cap classe gràfica.

La responsabilitat de cada capa és:

\begin{itemize}
    \item \textbf{Model}: conté les dades immutables del problema
          (\texttt{GameConfig}, \texttt{GameState}, \texttt{Keypad}) i els
          algorismes de programació dinàmica.
    \item \textbf{Vista}: implementada per \texttt{GameFrame}; dibuixa el
          teclat, els controls, els resultats i la reproducció pas a pas.
    \item \textbf{Controlador}: \texttt{GameController} llegeix el formulari,
          valida l'entrada, crea el teclat i l'algorisme seleccionat, i retorna
          un \texttt{GameResult} llest per mostrar.
\end{itemize}

La vista recalcula automàticament quan canvia qualsevol paràmetre, de manera que
no cal un botó de càlcul. Si la configuració és invàlida, el controlador captura
l'error de validació i envia un missatge a la vista en lloc d'executar cap algorisme.

El flux complet d'una modificació de la GUI és el següent:

\begin{enumerate}
    \item L'usuari canvia un camp, un selector o prem \textit{Nou aleatori}.
    \item \texttt{GameFrame} notifica \texttt{configurationChanged()} al seu
          \texttt{GameViewListener}.
    \item \texttt{GameController} llegeix un \texttt{GameInput} de la vista.
    \item El controlador construeix el \texttt{Keypad}, valida la configuració i
          crea un \texttt{GameConfig} amb índexs interns de jugador basats en 0.
    \item Segons el selector d'algorisme, crea \texttt{TopDownDPSolver} o
          \texttt{BottomUpDPSolver}.
    \item L'algorisme calcula el perdedor, la línia de reproducció i el nombre
          d'estats calculats.
    \item El controlador encapsula les dades en un \texttt{GameResult} i la vista
          redibuixa el teclat, els resultats i el pas actual de la reproducció.
\end{enumerate}

La Taula~\ref{tab:classes} resumeix tots els tipus definits a
\texttt{src/main/java}, no només el nucli algorísmic. Això deixa explícit com es
reparteix la responsabilitat entre el punt d'entrada, el model, el controlador,
la vista i el codi d'anàlisi experimental.

\begin{table*}[t]
\centering
\caption{Tots els tipus principals definits al projecte.}
\label{tab:classes}
\scriptsize
\begin{tabular}{@{}lll p{0.56\textwidth}@{}}
\toprule
\textbf{Paquet} & \textbf{Tipus} & \textbf{Nom} & \textbf{Responsabilitat} \\
\midrule
arrel & classe & \texttt{Main} & Punt d'entrada de la GUI; crea la vista i el controlador. \\
arrel & classe & \texttt{BenchmarkMain} & Punt d'entrada de consola per executar el benchmark i imprimir CSV. \\
model & classe & \texttt{Keypad} & Teclat immutable, disposició de valors i moviments vàlids precomputats. \\
model & enumeració & \texttt{KeypadMode} & Distingeix entre teclat estàndard i teclat aleatori. \\
model & enumeració & \texttt{SolverMode} & Selecciona l'algorisme \textit{top-down} o \textit{bottom-up}. \\
model & \textit{record} & \texttt{GameInput} & Configuració llegida de la vista abans de validar-la. \\
model & \textit{record} & \texttt{GameConfig} & Configuració validada i preparada per ser usada pels algorismes. \\
model & \textit{record} & \texttt{GameState} & Estat del joc $(t,l,j)$ emprat com a clau conceptual de la DP. \\
model & \textit{record} & \texttt{GameResult} & Resultat preparat per a la vista: perdedor, moviments inicials, estatística i reproducció. \\
model & \textit{record} & \texttt{ReplayStep} & Informació d'un pas concret de la reproducció òptima. \\
model & classe abstracta & \texttt{Solver} & Contracte comú dels algorismes i lògica compartida de reproducció. \\
model & classe & \texttt{TopDownDPSolver} & Implementació recursiva amb memoització a demanda. \\
model & classe & \texttt{BottomUpDPSolver} & Implementació iterativa que omple tota la taula d'estats. \\
controller & classe & \texttt{GameController} & Llegeix la configuració de la vista, valida, crea l'algorisme i retorna un \texttt{GameResult}. \\
view & interfície & \texttt{GameView} & Contracte que la vista exposa al controlador. \\
view & interfície & \texttt{GameViewListener} & Contracte d'esdeveniments que la vista emet cap al controlador. \\
view & classe & \texttt{GameFrame} & Finestra Swing amb formulari, teclat, resultats i navegació de la reproducció. \\
analysis & classe & \texttt{SolverBenchmark} & Executa sèries de benchmark i agrega mesures per configuració. \\
analysis & \textit{record} & \texttt{SolverMeasurement} & Una fila agregada de benchmark amb temps, memòria i estats calculats. \\
\bottomrule
\end{tabular}
\end{table*}

\subsection{Entrada, Validació i Conversió}

La vista treballa amb \texttt{GameInput}, un tipus \textit{record} de Java pensat
per representar el formulari tal com el veu l'usuari. En aquest objecte, el
jugador inicial usa indexació que comença a 1, perquè a la interfície és natural
mostrar ``Jugador 1'', ``Jugador 2'', etc. En canvi, els algorismes treballen amb
indexació que comença a 0, ja que això simplifica el càlcul del següent torn amb
l'operació mòdul. Aquesta conversió es fa només
al controlador, en el moment de crear \texttt{GameConfig}.

La validació està repartida en dos nivells. El controlador comprova les regles
que depenen de la interacció de diversos camps: que el límit sigui superior al
total inicial, que la darrera tecla sigui 0 o pertanyi al teclat, i que el
jugador inicial estigui dins el rang de jugadors. Després, \texttt{GameConfig}
repeteix les invariants essencials del model. Aquesta duplicació parcial és
intencionada: el controlador dona missatges d'error pensats per l'usuari, mentre
que el model es protegeix de qualsevol ús incorrecte des del benchmark o altres crides internes.

També hi ha una diferència entre \texttt{GameInput} i \texttt{GameResult}.
\texttt{GameInput} és entrada bruta de la vista; \texttt{GameResult} és sortida
preparada per mostrar. Inclou el teclat materialitzat, els moviments vàlids de
l'estat inicial, el jugador perdedor, el guanyador si només hi ha dos jugadors,
el nombre d'estats calculats i la llista de passos de reproducció. La vista no ha de
deduir cap regla del joc a partir d'aquest resultat; simplement el presenta.

\subsection{Representació dels Estats}

La classe \texttt{GameState} és un tipus \textit{record} immutable amb tres camps: total actual,
darrera tecla i torn. Aquesta immutabilitat és important per a la versió
\textit{top-down}, perquè els estats s'usen com a claus dins un \texttt{HashMap}.
Si un estat fos mutable, canviar-ne un camp després d'inserir-lo al mapa podria
trencar la correspondència entre \texttt{hashCode}, \texttt{equals} i la posició
interna de la taula hash.

La implementació evita guardar estats terminals amb $t\geq L$. Quan un moviment
arriba o supera el límit, el perdedor és conegut immediatament: és el jugador
que acaba de moure. Això redueix la taula a totals $0,\ldots,L-1$ i simplifica
el \textit{bottom-up}, perquè totes les posicions guardades representen estats on
encara s'ha de jugar.

El camp \texttt{lastPlayed} pot valer 0. Aquest valor no és una tecla real, sinó
un sentinella que representa el primer moviment. Per això la taula bottom-up usa
la dimensió \texttt{keyCount + 1}: les posicions de 1 a $K$ corresponen a tecles,
i la posició 0 correspon a l'estat inicial sense tecla prèvia.

\subsection{Model del Teclat}

\texttt{Keypad} és immutable. El constructor valida que les dimensions estiguin
dins el rang permès, que hi hagi exactament $W\cdot H$ valors i que aquests
formin una permutació sense duplicats. Durant la construcció es precomputen els
moviments vàlids de cada tecla. Això fa que els algorismes puguin consultar
\texttt{moves(lastPlayed)} sense tornar a recórrer la graella.

El teclat estàndard es construeix per files de dalt a baix amb els valors més
alts a la part superior. El teclat aleatori utilitza una permutació de les mateixes
tecles; a la GUI, el botó \textit{Nou aleatori} canvia la llavor i força una nova
disposició.

Internament la disposició es guarda en un vector unidimensional en ordre per files.
La posició $(r,c)$ es transforma en l'índex $r\cdot W+c$. Aquesta representació
és compacta i evita crear objectes per a cada casella. Per calcular els moviments
d'una tecla, el constructor localitza la seva fila i columna i copia tots els
valors de la mateixa fila i de la mateixa columna, excepte la pròpia tecla. El
resultat es guarda ordenat. L'ordre ascendent és rellevant perquè els algorismes
usen el primer moviment segur com a criteri de desempat i, a més, permet aturar
el recorregut quan apareix el primer moviment terminal: si un valor ja fa que
$t+m\geq L$, tots els moviments següents també seran terminals.

El cost de precomputar moviments és baix. Hi ha $K$ tecles i per a cada una es
recorren $W-1$ posicions de fila i $H-1$ posicions de columna. Per tant, el cost
de preparació del teclat és $O(K(W+H))$, molt menor que el cost potencial de
resoldre tots els estats quan $L$ o $P$ creixen. A canvi, cada consulta
\texttt{moves(lastPlayed)} pot retornar una còpia defensiva ja preparada. Les
còpies defensives impedeixen que una classe externa modifiqui el vector intern i
alteri les regles del joc després de construir el teclat.

\subsection{Algorismes}

Els dos algorismes hereten de \texttt{Solver}, que defineix el contracte comú:
calcular el perdedor, consultar el perdedor d'un estat arbitrari, indicar quants
estats s'han calculat i construir una reproducció òptima. El mètode de
reproducció també és comú, de manera que les dues implementacions comparteixen
la mateixa política de desempat.

\texttt{TopDownDPSolver} implementa la recurrència de l'Eq.~\ref{eq:recurrence}
amb un mapa \texttt{HashMap}. Quan rep una consulta, primer mira si
l'estat ja és a la memòria; si no hi és, avalua els moviments vàlids i guarda el
perdedor calculat. Aquesta variant és especialment útil quan l'estat inicial no
necessita visitar tota la taula teòrica.

El mètode central de la versió \textit{top-down} és \texttt{loser(state)}. La seva
estructura reflecteix directament la recurrència: per defecte assumeix que perd
el jugador actual; després revisa els moviments vàlids en ordre ascendent. Els
moviments terminals no es desenvolupen perquè ja fan perdre el jugador actual, i
com que els moviments estan ordenats, en trobar el primer moviment terminal es pot
aturar tot el bucle.
Per a cada moviment segur es consulta recursivament el successor. Si el perdedor
del successor no és el jugador actual, s'ha trobat una jugada que evita la derrota
pròpia i es pot aturar el bucle. Si cap moviment compleix aquesta condició,
l'assumpció inicial queda confirmada.

\begin{algorithm}[H]
\caption{Esquema de l'algorisme top-down.}
\label{alg:topdown}
\begin{algorithmic}[1]
\STATE \textbf{function} \textsc{Loser}$(t,l,j)$
\IF{estat és a \textit{memo}}
    \RETURN valor memoitzat
\ENDIF
\STATE $perdedor \leftarrow j$
\FORALL{$m \in M(l)$ en ordre ascendent}
    \IF{$t+m \geq L$}
        \STATE \textbf{break}
    \ENDIF
    \STATE $p \leftarrow \textsc{Loser}(t+m,m,(j+1)\bmod P)$
    \IF{$p \neq j$}
        \STATE $perdedor \leftarrow p$; \textbf{break}
    \ENDIF
\ENDFOR
\STATE guarda $perdedor$ a \textit{memo}
\RETURN $perdedor$
\end{algorithmic}
\end{algorithm}

La recursió no pot entrar en cicles perquè el total augmenta estrictament. La
profunditat màxima està limitada pel nombre de jugades abans d'arribar a $L$; en
el pitjor cas, si sempre es prem la tecla 1, és com a màxim $L-t_0$. En les
execucions realitzades això no ha creat problemes de pila; a més, la memoització
evita que branques diferents recalculin els mateixos successors.

\texttt{BottomUpDPSolver} crea dues taules tridimensionals de enters: una per al
perdedor de cada estat i una per al moviment triat. Les dimensions són
$L \cdot (K + 1) \cdot P$. El bucle principal recorre els
totals de \texttt{limit - 1} fins a \texttt{0}; com que qualsevol successor segur
té total més gran, el valor ja està disponible a la taula.

La taula \texttt{losers} guarda el jugador perdedor de cada estat. La taula
\texttt{moves} guarda el moviment escollit per a la reproducció. Aquesta segona
taula no és imprescindible per respondre qui perd, però evita tornar a recórrer
els moviments quan la GUI demana la línia òptima. El cost extra és una altra
taula d'enters amb les mateixes dimensions, acceptable dins els límits de teclat
i límit de la pràctica.

\begin{algorithm}[H]
\caption{Esquema de l'algorisme bottom-up.}
\label{alg:bottomup}
\begin{algorithmic}[1]
\FOR{$t=L-1$ \textbf{downto} $0$}
    \FOR{$l=0$ \textbf{to} $K$}
        \FOR{$j=0$ \textbf{to} $P-1$}
            \STATE $perdedor \leftarrow j$
            \STATE $moviment \leftarrow$ primer element de $M(l)$
            \FORALL{$m \in M(l)$ en ordre ascendent}
                \IF{$t+m\geq L$}
                    \STATE \textbf{break}
                \ENDIF
                \IF{$losers[t+m][m][(j+1)\bmod P]\neq j$}
                    \STATE $perdedor \leftarrow losers[t+m][m][(j+1)\bmod P]$
                    \STATE $moviment \leftarrow m$; \textbf{break}
                \ENDIF
            \ENDFOR
            \STATE $losers[t][l][j] \leftarrow perdedor$
            \STATE $moves[t][l][j] \leftarrow moviment$
        \ENDFOR
    \ENDFOR
\ENDFOR
\end{algorithmic}
\end{algorithm}

La part crítica és l'ordre del primer bucle. Si es recorreguessin els totals de
0 cap amunt, un estat podria necessitar un successor que encara no està calculat.
Recorrent de $L-1$ cap a 0, qualsevol successor segur $t+m$ és més gran que $t$ i
ja ha estat emplenat. Els successors terminals no necessiten entrada a la taula,
perquè la seva resposta és directa. Igualment, el recorregut dels moviments es pot
aturar al primer successor terminal perquè els valors legals es processen en ordre
ascendent.

\subsection{Reproducció de la Partida Òptima}

La GUI no només ha de mostrar el perdedor, sinó també una possible partida
òptima. Aquesta funcionalitat es concentra a \texttt{Solver.replay()}, compartida
pels dos algorismes. El mètode comença amb l'estat inicial i afegeix un
\texttt{ReplayStep} sense moviment, que representa la situació de partida. Després
repeteix el procés següent: demanar \texttt{chosenMove(state)}, sumar el moviment,
crear un nou pas de reproducció i, si el total encara és inferior al límit, avançar al
següent torn.

Cada \texttt{ReplayStep} conté les dades que necessita la vista: número de pas,
total anterior, total nou, darrera tecla, jugador que ha mogut, moviment fet,
següent jugador, moviments vàlids del pas següent i indicador de terminalitat.
Això evita que \texttt{GameFrame} hagi de reconstruir informació del model. Per
exemple, quan un pas és terminal, la llista de moviments vàlids següents és buida
i la vista només ha d'indicar que la partida ha acabat.

El desempat també és important en aquesta part. Si hi ha diversos moviments que
eviten perdre, es tria el primer en ordre ascendent. Si no n'hi ha cap, tots els
moviments són perdedors per al jugador actual; en aquest cas \texttt{chosenMove}
retorna el primer moviment vàlid com a línia representativa. Això dona una
partida reproduïble i fa que es pugui comparar la reproducció sencera entre
\textit{top-down} i \textit{bottom-up}.

\subsection{Interfície Gràfica}

La interfície Swing \cite{ref_gfg_swing} està organitzada en panells de configuració,
teclat, resultats i reproducció. L'usuari pot modificar dimensions, total inicial,
límit, darrera tecla, nombre de jugadors, jugador inicial, mode de teclat i mode
d'algorisme. El teclat es mostra com una calculadora real: les tecles vàlides de
l'estat actual es ressalten en blau, la darrera tecla jugada en groc i les tecles
invàlides en gris.

El panell de resultats mostra el jugador que perd, el guanyador quan $P=2$, els
moviments vàlids inicials i el nombre d'estats calculats. La reproducció permet
avançar i retrocedir pels passos de la partida òptima. A cada pas es veu qui ha
mogut, quina tecla ha premut, quin total s'ha obtingut i quin jugador mou a
continuació. Quan el pas és terminal, la vista indica que aquell moviment ha
arribat o superat el límit.

La interfície està pensada perquè el resultat sigui visible sense passos manuals
innecessaris. Els camps numèrics utilitzen components de selecció amb límits
raonables, els desplegables controlen el mode de teclat i el mode d'algorisme, i
qualsevol canvi llança una nova resolució. Això permet comparar ràpidament la
mateixa partida amb \textit{top-down} i \textit{bottom-up}, o veure com canvia el
perdedor en modificar el límit.

El dibuix del teclat usa el resultat del controlador, no recalcula moviments. La
vista rep el vector de moviments vàlids i el pas de reproducció actual. Amb aquesta
informació assigna colors: blau per a tecles vàlides, groc per a la darrera tecla
del pas i gris per a tecles no disponibles. Aquesta codificació visual ajuda a
comprovar manualment que la restricció de fila i columna s'està aplicant bé.

Per a partides amb més de dos jugadors, el camp \textit{Guanya} mostra un guió.
Aquesta decisió evita donar una resposta enganyosa. En un joc de $P>2$, saber qui
perd no determina un únic guanyador; simplement determina quin jugador queda
eliminat per la seqüència òptima considerada. El camp \textit{Perd} sí que es
manté per a qualsevol $P$.

\subsection{Gestió d'Errors i Casos Límit}

Els errors d'entrada es gestionen sense interrompre l'aplicació. Si el límit no
és superior al total inicial, si la darrera tecla no pertany al teclat o si el
jugador inicial és invàlid, el controlador no crea cap algorisme i la vista mostra
el missatge corresponent. Això és important perquè el recàlcul és automàtic: amb
un formulari interactiu és normal que l'usuari passi temporalment per valors
invàlids mentre està editant.

També s'han tractat dos casos especials del joc. El primer és el moviment inicial
amb \texttt{lastPlayed = 0}; en aquest cas \texttt{Keypad.moves(0)} retorna totes
les tecles ordenades. El segon és l'estat on tots els moviments vàlids són
terminals o condueixen a perdre: l'algorisme ha de retornar el jugador actual com a
perdedor, però encara ha de poder construir una reproducció. Per això es guarda
un moviment de reserva, el primer vàlid, que només s'usa quan no hi ha cap
moviment salvador.

\subsection{Benchmark}

Per separar les mesures de rendiment de la GUI s'ha creat el paquet
\texttt{analysis}.

\texttt{SolverBenchmark} executa els dos algorismes sobre una
llista de configuracions, aplica escalfaments previs i després calcula mitjanes
de les repeticions mesurades. \texttt{BenchmarkMain} imprimeix els resultats en
CSV, cosa que permet copiar-los directament a la memòria o a una eina externa.

La comanda utilitzada per a les mesures de la Secció~\ref{sec:resultats} és:

\begin{lstlisting}[language=bash,caption={Execució del benchmark de P5.},label={lst:benchmark}]
mvn -q exec:java \
  -Dexec.mainClass="com.serafinebot.p5.BenchmarkMain" \
  -Dexec.args="3 10"
\end{lstlisting}

El benchmark mesura la materialització del teclat i de la configuració, la
construcció de l'algorisme, el càlcul del perdedor i la generació de la
reproducció. No mesura el dibuix de la interfície Swing. El temps de CPU s'obté
amb \texttt{ThreadMXBean}; la memòria és una diferència
aproximada de memòria dinàmica abans i després de cada execució, útil per observar tendències.

Les configuracions per defecte del benchmark combinen dues sèries. La primera
recorre teclats quadrats de $2\times2$ fins a $5\times5$ amb $L=31$; la segona
fixa el teclat a $3\times3$ i fa créixer el límit fins a $10000$. Totes parteixen
de total 0, sense darrera tecla, dos jugadors i jugador inicial 1. S'ha triat
aquest conjunt perquè creix en dues direccions diferents: augmentar el límit fa
créixer la dimensió de totals, mentre que augmentar el teclat fa créixer tant el
nombre de valors de \texttt{lastPlayed} com la ramificació del primer moviment.

Cada fila de CSV inclou el nombre d'estats de la taula completa i el nombre
d'estats realment calculats per l'algorisme. En \textit{bottom-up} aquests dos valors
coincideixen sempre. En \textit{top-down}, el segon valor és la mida del mapa de
memoització després de resoldre i generar la reproducció. Aquesta diferència és la
magnitud més important per interpretar els resultats, perquè és independent del
soroll temporal propi d'executar microexecucions de rendiment sobre la JVM. A més,
el benchmark no inclou el cost de pintar la GUI ni la interacció de l'usuari;
només mesura la preparació mínima de l'entrada, la part algorísmica i la
construcció de la reproducció.

%% ================================================================
\section{Resultats i Discussió}
\label{sec:resultats}

\subsection{Protocol Experimental}

Les mesures s'han pres amb la comanda del Llistat~\ref{lst:benchmark}. Per a
cada configuració s'han executat 3 escalfaments i 10 repeticions mesurades. Els
teclats són estàndard, el total inicial és $0$, la darrera tecla és $0$, el
jugador inicial és el jugador 1 i hi ha 2 jugadors. S'han variat dues dimensions
de l'entrada: la mida del teclat i el límit de pèrdua. Per estudiar el límit
s'ha usat una sèrie llarga
$L\in\{31,60,100,200,500,1000,2000,5000,10000\}$ amb teclat $3\times3$.

Com que les instàncies són petites, els temps absoluts estan a l'escala de
microsegons i poden tenir soroll de JVM, planificador i recol·lector de brossa.
Per això la discussió dona més pes al nombre d'estats calculats, que és una
mesura directa de la feina estructural feta per cada algorisme. Els temps de CPU i
la variació de memòria dinàmica es mantenen com a indicadors pràctics complementaris.

\subsection{Creixement amb el Límit}

La Taula~\ref{tab:limit} fixa el teclat a $3\times3$ i mostra què passa quan
creix el límit $L$. La taula completa de l'algorisme \textit{bottom-up} creix de
forma exactament lineal amb $L$, tal com prediu l'Eq.~\ref{eq:states}. En canvi,
la versió \textit{top-down} visita menys estats, però també creix quan el límit
permet partides més llargues.

\begin{table}[H]
\centering
\caption{Sèrie de límits amb teclat $3\times3$ i 2 jugadors.}
\label{tab:limit}
\scriptsize
\begin{tabular}{@{}rrrrrr@{}}
\toprule
\textbf{$L$} & \textbf{Taula} & \textbf{TD est.} & \textbf{TD/taula} & \textbf{TD CPU ($\mu$s)} & \textbf{BU CPU ($\mu$s)} \\
\midrule
31     & 620     & 314    & 50,6\% & 189,2   & 120,8 \\
60     & 1.200   & 743    & 61,9\% & 89,4    & 64,9 \\
100    & 2.000   & 1.408  & 70,4\% & 162,7   & 94,6 \\
200    & 4.000   & 2.438  & 61,0\% & 270,0   & 265,0 \\
500    & 10.000  & 5.438  & 54,4\% & 546,3   & 460,5 \\
1.000  & 20.000  & 10.438 & 52,2\% & 605,8   & 766,0 \\
2.000  & 40.000  & 20.438 & 51,1\% & 1.485,8 & 1.486,0 \\
5.000  & 100.000 & 50.438 & 50,4\% & 5.246,8 & 4.015,3 \\
10.000 & 200.000 & 100.438 & 50,2\% & 10.317,1 & 7.934,6 \\
\bottomrule
\end{tabular}
\end{table}

\begin{figure}[H]
\centering
\begin{tikzpicture}
\begin{axis}[
    width=\columnwidth,
    height=4.7cm,
    xlabel={Límit $L$},
    ylabel={Estats calculats},
    xmode=log,
    log basis x={10},
    ymin=0,
    grid=major,
    legend pos=north west,
    legend style={font=\scriptsize, fill=white, fill opacity=0.85, draw opacity=1},
    tick label style={font=\scriptsize},
    label style={font=\scriptsize}
]
\addplot+[mark=*] coordinates {(31,314) (60,743) (100,1408) (200,2438) (500,5438) (1000,10438) (2000,20438) (5000,50438) (10000,100438)};
\addlegendentry{Top-down}
\addplot+[mark=square*] coordinates {(31,620) (60,1200) (100,2000) (200,4000) (500,10000) (1000,20000) (2000,40000) (5000,100000) (10000,200000)};
\addlegendentry{Bottom-up}
\end{axis}
\end{tikzpicture}
\caption{Creixement dels estats calculats amb teclat $3\times3$.}
\label{fig:limit_states}
\end{figure}

El resultat més rellevant és que \textit{bottom-up} calcula sempre el 100\% dels
estats de la taula. \textit{Top-down} en calcula una fracció menor: en aquesta
sèrie comença entre el 50\% i el 70\% i, per als límits grans, s'estabilitza prop
del 50\%. Això mostra que el benefici de la memoització a demanda no és només un
efecte de límits petits. La Fig.~\ref{fig:limit_states} mostra que els dos
algorismes creixen amb $L$, però \textit{bottom-up} segueix exactament la taula
completa, mentre que \textit{top-down} queda sistemàticament per davall.

\subsection{Creixement amb la Mida del Teclat}

La Taula~\ref{tab:size} fixa $L=31$ i mostra el creixement quan augmenta el
nombre de tecles. El nombre teòric d'estats creix amb $K+1$, i el factor de
ramificació també creix perquè al primer moviment totes les tecles són vàlides.

\begin{table}[H]
\centering
\caption{Estats calculats quan creix el teclat, amb $L=31$.}
\label{tab:size}
\scriptsize
\begin{tabular}{@{}lrrrrr@{}}
\toprule
\textbf{Teclat} & \textbf{$K$} & \textbf{$B$} & \textbf{Taula} & \textbf{TD estats} & \textbf{TD/taula} \\
\midrule
$2\times2$ & 4  & 4  & 310  & 62  & 20,0\% \\
$3\times3$ & 9  & 9  & 620  & 314 & 50,6\% \\
$4\times4$ & 16 & 16 & 1054 & 416 & 39,5\% \\
$5\times5$ & 25 & 25 & 1612 & 563 & 34,9\% \\
\bottomrule
\end{tabular}
\end{table}

\begin{figure}[H]
\centering
\begin{tikzpicture}
\begin{axis}[
    ybar,
    width=\columnwidth,
    height=4.8cm,
    ylabel={Top-down / taula (\%)},
    symbolic x coords={$2\times2$,$3\times3$,$4\times4$,$5\times5$},
    xtick=data,
    ymin=0,
    ymax=60,
    bar width=10pt,
    grid=major,
    tick label style={font=\scriptsize},
    label style={font=\scriptsize},
    x tick label style={font=\scriptsize}
]
\addplot coordinates {($2\times2$,20.0) ($3\times3$,50.6) ($4\times4$,39.5) ($5\times5$,34.9)};
\end{axis}
\end{tikzpicture}
\caption{Percentatge de la taula completa visitat per \textit{top-down} amb $L=31$.}
\label{fig:size_ratio}
\end{figure}

L'algorisme \textit{bottom-up} coincideix sempre amb la columna \textit{Taula},
perquè no filtra estats segons l'assolibilitat des de l'estat inicial. La versió
\textit{top-down}, en canvi, queda entre el 20\% i el 50,6\% en aquests casos. El
percentatge no és monòton perquè depèn de les sumes concretes que es poden
aconseguir amb cada teclat i de la política de desempat, com es veu a la
Fig.~\ref{fig:size_ratio}. Tot i així, confirma que
la memoització evita una part important de la taula completa.

\subsection{Temps de CPU i Memòria}

La Fig.~\ref{fig:cpu_3x3} representa els temps de CPU de la mateixa sèrie llarga
de límits de la Taula~\ref{tab:limit}. Els temps absoluts continuen essent petits,
però amb nou valors de $L$ ja es veu una tendència creixent. També es veu que el
temps no segueix exactament el nombre d'estats: les constants de la JVM, el cost
del \texttt{HashMap} i la localitat dels vectors primitius fan que hi hagi canvis
de lideratge entre les dues variants.

\begin{figure}[H]
\centering
\begin{tikzpicture}
\begin{axis}[
    ybar,
    width=\columnwidth,
    height=5.0cm,
    ylabel={CPU ($\mu$s)},
    xlabel={Límit $L$},
    ymode=log,
    log basis y={10},
    symbolic x coords={31,60,100,200,500,1000,2000,5000,10000},
    xtick=data,
    bar width=3pt,
    enlarge x limits=0.08,
    legend pos=north west,
    legend style={font=\scriptsize, fill=white, fill opacity=0.85, draw opacity=1},
    grid=major,
    tick label style={font=\scriptsize},
    label style={font=\scriptsize},
    x tick label style={font=\scriptsize, rotate=45, anchor=east}
]
\addplot+[bar shift=-2pt] coordinates {(31,189.2) (60,89.4) (100,162.7) (200,270.0) (500,546.3) (1000,605.8) (2000,1485.8) (5000,5246.8) (10000,10317.1)};
\addlegendentry{Top-down}
\addplot+[bar shift=2pt] coordinates {(31,120.8) (60,64.9) (100,94.6) (200,265.0) (500,460.5) (1000,766.0) (2000,1486.0) (5000,4015.3) (10000,7934.6)};
\addlegendentry{Bottom-up}
\end{axis}
\end{tikzpicture}
\caption{Temps de CPU amb teclat $3\times3$ per a la sèrie llarga de límits.}
\label{fig:cpu_3x3}
\end{figure}

En memòria, la mesura directa de la memòria dinàmica del procés no és prou estable
per representar-la com a gràfica: la JVM pot reutilitzar espai o activar el
recol·lector de brossa entre repeticions. Per això la comparació principal de
memòria és estructural. \textit{Bottom-up} reserva dues taules de mida
proporcional a $S_{max}$, mentre que \textit{top-down} guarda $R$ entrades dins un
\texttt{HashMap}. Així, \textit{top-down} sol tenir menys estats, però cada entrada
té més sobrecàrrega; \textit{bottom-up} guarda més posicions, però en vectors
primitius més compactes.

\subsection{Discussió General}

Les dues variants són correctes i comparteixen la mateixa política de joc òptim.
La diferència és quan i quant calculen. \textit{Bottom-up} és simple de raonar,
té accés $O(1)$ a qualsevol estat després de construir-se i el seu recompte
d'estats és totalment predictible. A canvi, fa feina sobre estats que potser mai
seran necessaris per a la configuració inicial.

\textit{Top-down} és més selectiu: només calcula allò que la recursió necessita.
És l'opció natural per a la GUI, perquè l'usuari normalment demana una sola
configuració inicial i no consultes arbitràries sobre tota la taula. El cost és
que cada estat memoitzat té més sobrecàrrega d'objectes, i que el temps depèn més
de la forma concreta de l'espai accessible.

Per tant, la selecció final no és absoluta. Si es volen moltes consultes sobre
la mateixa configuració global, \textit{bottom-up} pot compensar el cost inicial.
Si es vol resoldre una partida concreta i mostrar-ne una reproducció, la versió
\textit{top-down} sol fer menys feina.

%% ================================================================
\section{Conclusions}

La pràctica ha implementat una aplicació completa per resoldre el joc de la
calculadora amb programació dinàmica. El model del problema queda capturat per
l'estat $(t,l,j)$ i per una recurrència que retorna el jugador perdedor sota joc
òptim. La propietat que el total sempre augmenta permet evitar cicles i justifica
tant la memoització \textit{top-down} com la taula \textit{bottom-up}.

L'arquitectura MVC ha separat la lògica del joc de la interfície Swing. Això ha
facilitat validar el controlador sense obrir finestres i reutilitzar el mateix
model per al benchmark. La GUI permet configurar la partida, alternar entre
algorismes, usar teclats estàndard o aleatoris i seguir una línia òptima pas a pas.

Els resultats experimentals mostren que \textit{bottom-up} calcula sempre
$L\cdot(K+1)\cdot P$ estats, mentre que \textit{top-down} pot calcular-ne una
fracció menor. Els temps de CPU creixen en general amb la mida de la instància,
però no segueixen de manera directa el nombre d'estats: les constants
d'implementació poden fer que \textit{bottom-up} sigui més ràpid en alguns
límits grans. En memòria, \textit{top-down} redueix el nombre d'estats guardats,
però el \texttt{HashMap} introdueix més sobrecàrrega per estat que els vectors
primitius.

Com a treball futur es podrien afegir teclats editables manualment, comptar el
nombre total de línies de joc possibles independentment de la línia òptima i
ampliar el benchmark amb límits més grans o més jugadors per observar millor el
factor $P$ de la complexitat.

%% ================================================================
\begin{thebibliography}{9}

\bibitem{ref_gfg_dp}
GeeksforGeeks, ``Dynamic Programming or DP,'' 2026. [Online]. Available:
\url{https://www.geeksforgeeks.org/dynamic-programming/}

\bibitem{ref_mvc_wikipedia}
Wikipedia contributors, ``Model--view--controller,'' Wikipedia. [Online].
Available: \url{https://en.wikipedia.org/wiki/Model-view-controller}

\bibitem{ref_gfg_swing}
GeeksforGeeks, ``Introduction to Java Swing,'' 2025. [Online]. Available:
\url{https://www.geeksforgeeks.org/introduction-to-java-swing/}

\end{thebibliography}

\end{document}
